% TEMPLATE-MILC-v3.tex - plantilla maestra UNICA de las guias MILC v3.
%
% La usan las cuatro familias de guias, cada una con su driver:
%   · Grados 6-11          -> scripts/build-guias-g11.py
%   · Bebras               -> scripts/build-guias-bebras.py
%   · Semillero            -> scripts/build-guias-semillero.py
%   · Territorio interior  -> scripts/build-guias-territorio-interior.py
%
% Antes habia tres copias casi identicas (~400 lineas iguales, 6 distintas):
% cada mejora habia que aplicarla tres veces, y en la practica no se hacia
% (el motor de recursos vivia solo en dos de ellas). Lo que varia entre
% familias esta parametrizado en 6 placeholders que cada driver llena
% (se nombran aqui SIN su sintaxis real: el builder reemplaza por texto plano,
% tambien dentro de los comentarios, y un valor multilinea rompe el preambulo):
%
%   HEADER_IZQ        encabezado izquierdo de pagina
%   HEADER_DER        encabezado derecho de pagina
%   PORTADA_ETIQUETA  rotulo sobre el numero grande (GRADO/MOMENTO/MODULO)
%   PORTADA_SERIE     linea de serie de la portada
%   PORTADA_ID        identificador de la guia en la portada
%   PORTADA_CIRCULO   el circulito de grado (°), o vacio si no aplica
%   PDF_TITULO        titulo en texto plano (sin \textbf ni comillas TeX) para
%                     los metadatos del PDF (/Title); lo produce lib_xelatex.texto_pdf
%
% Composicion (auditoria 2026-09-03): las bandas de estacion piden espacio
% (\Needspace*) y prohiben el salto justo despues (\nopagebreak), asi no quedan
% huerfanas al pie; las cajas largas (softbox, drawbox, infoband, puentebox) son
% `breakable`, asi una caja de media pagina ya no arrastra todo a la siguiente
% dejando la anterior al 40 %. Todo texto sobre mostaza o verde va en milcNegro
% (blanco sobre #E5B400 da 1,93:1 y sobre #58A923 2,95:1; ninguno pasa WCAG AA).
% El resto de claves las documenta content/guias/_SCHEMA.md. El script valida
% que no queden placeholders sin reemplazar antes de escribir el archivo final.

% Etiquetado PDF (latex-lab): /Lang, estructura y texto alternativo de las
% figuras, para lectores de pantalla. Debe ir ANTES de \documentclass.
\DocumentMetadata{lang=es-CO,tagging=on}
\documentclass[11pt,letterpaper]{article}
% headheight=24pt: la cabecera derecha lleva dos lineas (identidad de la guia
% y estacion actual). headsep se acorta en la misma medida para que el bloque
% de texto quede exactamente donde estaba con headheight=12pt.
\usepackage[margin=2.0cm,headheight=24pt,headsep=13pt]{geometry}
\usepackage[table]{xcolor}
\usepackage{fontspec}
\usepackage[spanish]{babel}
\usepackage{tikz}
% Librerías TikZ para los diagramas de las guías (bloque `recursos:`):
%  · babel  → babel en español vuelve activos caracteres como «>», y TikZ los usa
%             en sus opciones (>=stealth, ->). Sin esta librería, cualquier
%             diagrama con flechas revienta con «\language@active@arg> has an
%             extra }». Debe ir ANTES de que se dibuje nada.
%  · positioning        → sintaxis «below left=2pt and 3pt».
%  · decorations.pathreplacing → llaves ({brace}) para señalar rangos.
\usetikzlibrary{babel,positioning,decorations.pathreplacing}
\usepackage{graphicx}
% Circuitos: el trabajo de electrónica se hace hoy con micro:bit y sus sensores
% integrados (MakeCode, sin cableado), así que no se necesitan esquemáticos.
% Si algún día entran montajes con protoboard, basta descomentar esta línea:
% los diagramas .tex/.tikz declarados en `recursos:` ya se incrustan solos.
% \usepackage{circuitikz}
\usepackage[most]{tcolorbox}
\usepackage{tabularx}
\usepackage{array}
\usepackage{fancyhdr}
\usepackage{titlesec}
\usepackage[document]{ragged2e}  % bandera a la derecha en todo el cuerpo
\usepackage{enumitem}
\usepackage{microtype}
\usepackage{etoolbox}
\usepackage{needspace}
% hyperref al final del preambulo: metadatos del PDF (titulo, autor, idioma,
% familia), marcadores por estacion (\pdfbookmark) y enlaces sin recuadro.
\usepackage[hidelinks,
  pdftitle={Mi correo institucional — la primera dirección que firmas como tuya},
  pdfauthor={PhD. Álvaro Cárdenas Orozco},
  pdfsubject={Tecnología e Informática · Grado 6 · Guía 3 · MILC v3},
  pdflang={es-CO},
  bookmarksopen=true,bookmarksnumbered=false]{hyperref}

% Helvetica Neue no trae la flecha U+2192, asi que cada «→» escrito literalmente
% en el YAML se compilaba a NADA, en silencio (462 flechas en 61 guias: rutas del
% tipo «paleta → tipografia → lockup» salian rotas). Mapeamos el caracter a su
% version matematica, que si tiene glifo. Asi el autor puede seguir escribiendo
% «→» comodamente en el YAML.
\usepackage{newunicodechar}
\newunicodechar{→}{\ensuremath{\rightarrow}}
% Mismo problema con los simbolos de marca y vinetas: el «✓» de «una palomita ✓»
% salia como cuadrito vacio en el PDF de todas las guias que lo usaban.
\usepackage{pifont}
\newunicodechar{✓}{\ding{51}}
\newunicodechar{✗}{\ding{55}}
\newunicodechar{✔}{\ding{52}}
\newunicodechar{•}{\textbullet}
\newunicodechar{≥}{\ensuremath{\geq}}
\newunicodechar{≤}{\ensuremath{\leq}}

\setmainfont{Helvetica Neue}
\setsansfont{Helvetica Neue}
\setlength{\parindent}{0pt}
\setlength{\parskip}{6pt}
% Sin particion de palabras: en 6.º-7.º una palabra cortada por guion al final
% de linea («Comuni-cacion») frena la lectura mucho mas que una linea corta.
\hyphenpenalty=10000
\exhyphenpenalty=10000
\pretolerance=2000
\tolerance=3000
\setlength{\emergencystretch}{3em}
\linespread{1.10}
\renewcommand{\arraystretch}{1.25}
\setlist{leftmargin=5mm,itemsep=1mm,topsep=1mm}
\newcolumntype{Y}{>{\RaggedRight\arraybackslash}X}

\definecolor{milcMagenta}{HTML}{E6007E}
\definecolor{milcVino}{HTML}{5A0038}
\definecolor{milcTurquesa}{HTML}{0093A5}
\definecolor{milcVerde}{HTML}{58A923}
\definecolor{milcMostaza}{HTML}{E5B400}
\definecolor{milcOcre}{HTML}{B86600}
\definecolor{milcNegro}{HTML}{241816}
\definecolor{milcGris}{HTML}{F7F7F4}
\definecolor{milcLinea}{HTML}{E8E2DD}
\definecolor{milcGradePrimary}{HTML}{0066FF}
\definecolor{milcGradeDark}{HTML}{003D99}
\definecolor{milcGradeSoft}{HTML}{D6E8FF}
\definecolor{milcGradeAccent}{HTML}{CFE7FF}
\definecolor{milcGradeText}{HTML}{FFFFFF}
\color{milcNegro}

\pagestyle{fancy}
\fancyhf{}
% Cabecera de dos lineas: identidad de la guia arriba y, a la derecha y en
% gris, la estacion en curso (\markright desde \milcestacion). Antes de la
% primera estacion la segunda linea va vacia; el \strut mantiene la altura
% para que la linea de arriba no baile entre paginas.
\lhead{\small Tecnología e Informática · Grado 6\\[-1pt]{\footnotesize\strut}}
\rhead{\small Guía 3 · MILC v3\\[-1pt]{\footnotesize\strut\color{milcNegro!65}\rightmark}}
\cfoot{\small\thepage}
\renewcommand{\headrulewidth}{0pt}

% Sobre mostaza (#E5B400) y verde (#58A923) el texto va en milcNegro; sobre el
% resto de la paleta, en blanco. Se usa dentro de fonttitle/fontupper, donde un
% \color posterior manda sobre coltitle.
\newcommand{\milctintasobre}[1]{%
  \ifboolexpr{test{\ifstrequal{#1}{milcMostaza}} or test{\ifstrequal{#1}{milcVerde}}}%
    {\color{milcNegro}}{\color{white}}}

\newtcolorbox{titlebox}[1]{
  enhanced,colback=#1,colframe=#1,arc=7mm,boxrule=0pt,
  left=7mm,right=7mm,top=3mm,bottom=3mm,
  before skip=8pt,after skip=8pt,
  fontupper=\sffamily\bfseries\Large\color{white}
}

\newtcolorbox{softbox}[3]{
  enhanced,breakable,pad at break=3mm,
  colback=#2,colframe=#1,borderline west={4pt}{0pt}{#1},
  arc=2mm,boxrule=.4pt,left=4mm,right=4mm,top=3mm,bottom=3mm,
  before skip=6pt,after skip=8pt,title={#3},coltitle=white,
  colbacktitle=#1,fonttitle=\sffamily\bfseries\milctintasobre{#1},fontupper=\RaggedRight,
  attach boxed title to top left={xshift=3mm,yshift=-2mm},
  boxed title style={arc=2mm, boxrule=0pt}
}

\newtcolorbox{drawbox}[2]{
  enhanced,breakable,pad at break=4mm,
  colback=white,colframe=#1,arc=4mm,boxrule=1pt,
  left=5mm,right=5mm,top=4mm,bottom=4mm,before skip=8pt,after skip=8pt,
  title={#2},coltitle=white,colbacktitle=#1,fonttitle=\sffamily\bfseries\milctintasobre{#1},
  fontupper=\RaggedRight,
  attach boxed title to top left={xshift=4mm,yshift=-2mm},
  boxed title style={arc=3mm, boxrule=0pt}
}

\newtcolorbox{infoband}[1]{
  enhanced,breakable,pad at break=2mm,colback=#1!8,colframe=#1!50,arc=2mm,boxrule=.5pt,
  left=4mm,right=4mm,top=2mm,bottom=2mm,before skip=4pt,after skip=4pt,
  fontupper=\small\RaggedRight
}

% Puente narrativo entre secciones (contrato editorial Conéctate).
% Da continuidad: "Acabas de X. Ahora vas a Y porque Z."
% Visualmente: banda izquierda en color del grado, texto en itálica pequeña.
\newtcolorbox{puentebox}{
  enhanced,breakable,pad at break=2mm,colback=milcGradeSoft,colframe=milcGradeDark,
  borderline west={3pt}{0pt}{milcGradeDark},
  arc=0mm,boxrule=0pt,left=5mm,right=4mm,top=2.5mm,bottom=2.5mm,
  before skip=4pt,after skip=8pt,
  fontupper=\itshape\small\color{milcNegro}\RaggedRight
}
% La flecha va en modo matematico a proposito: Helvetica Neue NO tiene el glifo
% U+2192, asi que \textrightarrow se compilaba a NADA («Missing character: There
% is no → in font Helvetica Neue») y el puente salia sin flecha. En modo
% matematico la flecha viene de la fuente matematica y si se dibuja.
\newcommand{\puente}[1]{%
  \begin{puentebox}%
    \textcolor{milcGradeDark}{$\rightarrow$}\hspace{2mm}#1%
  \end{puentebox}%
}

\newcommand{\linead}[1]{\noindent\color{milcLinea}\rule{#1}{0.5pt}\color{milcNegro}}
\newcommand{\respuesta}[1]{\par\vspace{2mm}\foreach \n in {1,...,#1}{\noindent\color{milcLinea}\rule{\linewidth}{0.5pt}\par\vspace{5mm}}\color{milcNegro}}
\newcommand{\checkbox}{\textcolor{milcGradeDark}{\fbox{\phantom{x}}}\hspace{5pt}}

% ─── Listas aireadas para las guías socioemocionales ────────────────────
% Componer las verificaciones como «\checkbox texto\\» deja el texto que salta
% de línea debajo del cuadrito y apelmaza el bloque. Estos entornos alinean la
% continuación y airean los ítems. Los usa Territorio Interior; cualquier
% familia puede hacerlo.
\newlist{checklista}{itemize}{1}
\setlist[checklista]{
  label=\textcolor{milcGradeDark}{\fbox{\phantom{x}}},
  leftmargin=9mm, labelsep=3.2mm, itemsep=2.4mm,
  topsep=2.5mm, parsep=0pt, partopsep=0pt, before=\RaggedRight
}
\newlist{pasoslista}{enumerate}{1}
\setlist[pasoslista]{
  label=\textcolor{milcGradeDark}{\sffamily\bfseries\arabic*.},
  leftmargin=8mm, labelsep=3mm, itemsep=2.8mm,
  topsep=1mm, parsep=0pt, partopsep=0pt, before=\RaggedRight
}

\newcommand{\phasecard}[4]{
\begin{tcolorbox}[enhanced,colback=#3,colframe=#2,arc=3mm,boxrule=.5pt,height=3.05cm,valign=center,halign=center,left=1.5mm,right=1.5mm,top=2mm,bottom=2mm]
{\sffamily\bfseries\fontsize{22}{24}\selectfont\textcolor{#2}{#1}}\par
{\sffamily\bfseries\small #4}
\end{tcolorbox}
}

% ── Piezas del rediseno 2026 (legibilidad 6.º-11.º) ───────────────────
% Banda de estacion: reemplaza el titlebox macizo. Titulo a la izquierda,
% dato practico (tiempo/modalidad) a la derecha, en una sola linea.
% Nunca huerfana: pide 9 lineas libres (o salta de pagina rellenando con
% \vfill) y prohibe el salto justo despues de la banda. Nueve y no seis:
% la caja partible que sigue necesita ~80pt (titulo adosado, relleno y dos
% lineas) para arrancar en la misma pagina; con menos, tcolorbox la manda
% entera a la siguiente y la banda se queda sola. El penalty 10000 va
% en `after`, ANTES del espacio posterior: un `after skip` (aunque sea 0pt)
% es un \vskip tras la caja, y ese glue es un punto de corte legal que un
% \nopagebreak escrito despues de \end{tcolorbox} ya no cierra. Cada banda
% deja marcador PDF y actualiza la cabecera derecha.
\newcounter{milcestacion}
\newcommand{\milcestacion}[2]{%
\Needspace*{9\baselineskip}%
\stepcounter{milcestacion}%
\pdfbookmark[1]{#2}{milcestacion\themilcestacion}%
\markright{#2}%
\begin{tcolorbox}[enhanced,colback=#1,colframe=#1,arc=2mm,boxrule=0pt,
  left=5mm,right=5mm,top=2.6mm,bottom=2.6mm,before skip=11pt,
  after={\par\nopagebreak\vskip7pt}]
{\sffamily\bfseries\fontsize{15}{18}\selectfont\milctintasobre{#1}#2}
\end{tcolorbox}}

% Tarjeta de paso numerada: sustituye las tablas «Pilar / Aplicacion», que
% eran formato de consulta docente, no de estudiante. El numero va en un
% circulo sobre el borde para que la secuencia se lea de un vistazo.
\newcommand{\milcpaso}[3]{%
\Needspace{4\baselineskip}%
\begin{tcolorbox}[enhanced,colback=white,colframe=milcLinea,arc=1.5mm,boxrule=.9pt,
  left=14mm,right=4mm,top=3mm,bottom=3mm,before skip=4pt,after skip=4pt,
  fontupper=\RaggedRight,
  overlay={\node[anchor=north west,circle,fill=#2,text=white,inner sep=0pt,
    minimum size=7.4mm,font=\sffamily\bfseries\small]
    at ([xshift=3.6mm,yshift=-3.2mm]frame.north west) {#1};}]
#3
\end{tcolorbox}}

% Casilla marcable de tamano util con lapiz (la anterior era \fbox{\phantom{x}},
% de alto variable segun el texto que la seguia).
\newcommand{\milcmarca}{\framebox[3.8mm]{\rule{0pt}{3.4mm}}}

% Fila marcable de lista suelta (checks de la estacion 1).
\newcommand{\milccheckrow}[1]{%
\par\vspace{1.8mm}\noindent
\makebox[6.4mm][l]{\raisebox{-.55mm}{\textcolor{milcNegro}{\milcmarca}}}%
\parbox[t]{\dimexpr\linewidth-6.4mm\relax}{\RaggedRight #1}\par}

% Celda marcable dentro de la rubrica (tabla de autoevaluacion). Misma
% sangria francesa, pero pensada para vivir dentro de una columna X.
\newcommand{\milcopcion}[1]{%
\makebox[5.2mm][l]{\raisebox{-.55mm}{\textcolor{milcNegro}{\milcmarca}}}%
\parbox[t]{\dimexpr\linewidth-5.2mm\relax}{\RaggedRight #1}}

% ── Recursos visuales de la guía ──────────────────────────────────────
% Declarados en el bloque `recursos:` del YAML y emitidos por el builder.
% \guiaFigura   → imagen rasterizada o PDF (foto, carta celeste, montaje).
% \guiaDiagrama → fuente TikZ/circuitikz (.tex) incrustada como vector:
%                 es la vía para circuitos y esquemáticos editables.
% [#1] = texto alternativo (alt) · #2 = ruta relativa al .tex · #3 = pie
% El alt llega al PDF etiquetado (/Alt) y lo lee el lector de pantalla.
\newcommand{\guiaFigura}[3][]{%
\begin{center}
\includegraphics[width=0.88\linewidth,height=0.38\textheight,keepaspectratio,alt={#1}]{#2}\\[1.5mm]
{\footnotesize\itshape\textcolor{milcNegro!75}{#3}}
\end{center}
}

\newcommand{\guiaDiagrama}[2]{%
\begin{center}
\input{#1}\\[1.5mm]
{\footnotesize\itshape\textcolor{milcNegro!75}{#2}}
\end{center}
}

\begin{document}
\thispagestyle{empty}

% ── Portada ───────────────────────────────────────────────────────────
% Antes: un rectangulo de color a pagina completa. Se veia bien en pantalla,
% pero en la fotocopiadora del colegio se comia el toner y salia gris sucio.
% Ahora la banda ocupa el tercio superior y el resto va en blanco.
\begin{tikzpicture}[remember picture,overlay]
  \path[fill=milcGradePrimary]
    (current page.north west) rectangle ([yshift=-10.6cm]current page.north east);

  \node[anchor=north west,text=milcGradeText] at ([xshift=2.0cm,yshift=-1.55cm]current page.north west)
    {\sffamily\bfseries\fontsize{12}{14}\selectfont GRADO};
  \node[anchor=north west,text=milcGradeText,opacity=.85] at ([xshift=2.0cm,yshift=-2.35cm]current page.north west)
    {\sffamily\bfseries\fontsize{8.5}{10}\selectfont SERIE GUÍAS · TECNOLOGÍA E INFORMÁTICA};
  \node[anchor=north east,text=milcGradeText,opacity=.85,align=right] at ([xshift=-2.0cm,yshift=-1.55cm]current page.north east)
    {\sffamily\bfseries\fontsize{9}{12}\selectfont I.E. Sor María Juliana\\Cartago, Valle del Cauca};

  \node[anchor=west,text=milcGradeText] (gradonum) at ([xshift=1.85cm,yshift=-7.1cm]current page.north west)
    {\sffamily\bfseries\fontsize{112}{112}\selectfont 6};
  % El circulito de grado (°) solo aplica a las guias de grado («11°»); en
  % Bebras y Semillero el numero grande es un momento/modulo, no un grado.
  % Se posiciona relativo al nodo `gradonum`, no en coordenadas absolutas.
  \draw[milcGradeText,line width=1.6pt]
    ([xshift=2.0mm,yshift=-4.5mm]gradonum.north east) circle (.15cm);

  \node[anchor=west,text=milcGradeText,text width=11.4cm,align=left] at ([xshift=1.5cm,yshift=0.5cm]gradonum.east)
    {
     {\sffamily\bfseries\fontsize{9}{11}\selectfont\MakeUppercase{Período 1 · Comunicación e identidad digital}}\\[3mm]
     {\sffamily\bfseries\fontsize{21}{25}\selectfont\setlength{\baselineskip}{27pt}Mi correo\\[4pt]institucional ---\\[4pt]mi primera\\[4pt]dirección propia}\\[3.5mm]
     {\sffamily\bfseries\fontsize{15}{18}\selectfont Guía 3-6-TIC}
    };
\end{tikzpicture}

\vspace*{9.4cm}
\vfill

{\sffamily\bfseries\fontsize{8}{10}\selectfont\textcolor{milcGradeDark}{LO QUE TE LLEVAS HOY}}

\begin{tcolorbox}[enhanced,colback=white,colframe=milcNegro,arc=2mm,boxrule=1.6pt,
  left=5mm,right=5mm,top=4mm,bottom=4mm,before skip=3pt,after skip=12pt,
  fontupper=\RaggedRight\fontsize{11}{15.5}\selectfont]
\textbf{Correo institucional configurado y funcionando} (sabes el usuario, sabes la contraseña, sabes cómo abrirlo). Más \textbf{3 correos formales escritos en el cuaderno}: uno al profe de Tecnología, uno al coordinador del colegio, uno a un familiar lejano. Los 3 correos siguen las 5 partes que vas a aprender. Lo muestras antes de salir de la sala.
\end{tcolorbox}

\begin{tcolorbox}[enhanced,colback=milcGris,colframe=milcGris,arc=2mm,boxrule=0pt,
  left=5mm,right=5mm,top=3.5mm,bottom=3.5mm,after skip=12pt]
{\sffamily\bfseries\fontsize{8}{10}\selectfont\textcolor{milcNegro!55}{ESTA GUÍA ES DE}}\\[3mm]
\begin{tabularx}{\linewidth}{X p{3.2cm} p{3.2cm}}
\linead{\linewidth} & \linead{3.0cm} & \linead{3.0cm}\\[-1mm]
{\fontsize{8.5}{10}\selectfont\textcolor{milcNegro!55}{Nombre completo}} &
{\fontsize{8.5}{10}\selectfont\textcolor{milcNegro!55}{Grupo}} &
{\fontsize{8.5}{10}\selectfont\textcolor{milcNegro!55}{Fecha}}\\
\end{tabularx}
\end{tcolorbox}

\vfill

\noindent\textcolor{milcLinea}{\rule{\linewidth}{1pt}}\\[2mm]
{\sffamily\bfseries\fontsize{9.5}{12}\selectfont Autor: Dr. Álvaro Cárdenas Orozco}\\[1mm]
{\sffamily\fontsize{8.5}{11}\selectfont\textcolor{milcNegro!55}{Metodología MILC · apertura ancestral · pensamiento computacional · triángulo Dussel–estoicismo–Floridi}}

\newpage

\section*{Apertura: Mi correo institucional --- la primera dirección que firmas como tuya}

\begin{softbox}{milcGradeDark}{milcGradeSoft}{Saber ancestral}
\textbf{Vuelve el pregonero, pero esta vez con cartas.} Antes de los celulares, las cartas viajaban semanas. Para que una carta llegara, tenía que llevar \textbf{una dirección clara}: \emph{``Señor Aurelio Martínez, casa de tapia blanca, frente al palo de mango, vereda La Plata, Cartago.''}. Si la dirección estaba mal escrita, la carta se perdía o llegaba tarde. El cartero, igual que el pregonero, conocía a las personas por sus direcciones: la dirección era casi parte del nombre. La gente cuidaba mucho cómo escribía esa dirección y a quién se la daba. \textbf{Una buena carta tenía cinco cosas:} a quién va, de parte de quién, sobre qué tema, el mensaje, y la firma. Hoy las cartas se llaman \textbf{correos electrónicos}, y siguen necesitando esas mismas cinco cosas para llegar bien.
\end{softbox}

\begin{softbox}{milcTurquesa}{milcGris}{Saber contemporáneo}
Tu \textbf{correo institucional} es tu primera dirección propia en internet. Termina en \texttt{@sormariajuliana.edu.co} (o el dominio del colegio). No la elegiste por gusto: el colegio te la asignó porque eres parte de la institución. Por eso es \textbf{seria}, no es como un correo de juego. La vas a usar para: recibir tareas del profe, enviar trabajos, comunicarte con coordinación, registrarte en plataformas educativas. Y --- esto es importante --- vas a aprender a \textbf{escribir un correo formal} con las 5 partes correctas. Esa habilidad la vas a usar toda la vida: para pedir un trabajo, para escribirle a un profesor de universidad, para hacer un reclamo en serio.
\end{softbox}

\begin{softbox}{milcMostaza}{milcGris}{Pregunta-puente}
\textit{Si tuvieras que escribirle al director del colegio para pedirle algo importante, ¿\textbf{cómo empezarías} el correo? ¿``Hola, qué más''? ¿``Buenas tardes''? ¿``Estimado señor director''? La respuesta importa.}
\end{softbox}

\begin{softbox}{milcVerde}{milcGris}{Saber y saber hacer}
Al terminar la clase: (1) podrás \textbf{identificar} las 5 partes de un correo formal; (2) sabrás \textbf{explicar} con tus palabras por qué cada parte importa; (3) habrás \textbf{configurado} tu correo institucional con contraseña segura; (4) habrás \textbf{escrito 3 correos formales} en el cuaderno.
\end{softbox}



\small
\begin{tabularx}{\textwidth}{>{\bfseries}p{3.0cm}Y}
\rowcolor{milcGradePrimary!12}Autor & Dr. Álvaro Cárdenas Orozco\\
DBA & Reconoce principios y conceptos propios de la tecnología, relacionando artefactos con su utilización segura (MEN, T\&I 6°)\\
Referentes & Saber ancestral del pregonero · Pensamiento computacional inicial · Cuidado de la palabra\\
Producto final & \textbf{Correo institucional configurado y funcionando} (sabes el usuario, sabes la contraseña, sabes cómo abrirlo). Más \textbf{3 correos formales escritos en el cuaderno}: uno al profe de Tecnología, uno al coordinador del colegio, uno a un familiar lejano. Los 3 correos siguen las 5 partes que vas a aprender. Lo muestras antes de salir de la sala.\\
\end{tabularx}
\normalsize

\puente{Hoy haces 4 cosas: lees un ejemplo de correo, aprendes las 5 partes, activas tu correo institucional, y escribes 3 correos formales.}

\milcestacion{milcGradeDark}{Ruta MILC: así vamos a aprender}
\begin{tabularx}{\textwidth}{>{\centering\arraybackslash}X>{\centering\arraybackslash}X>{\centering\arraybackslash}X>{\centering\arraybackslash}X}
\phasecard{1}{milcVerde}{milcGris}{Escucha\\[-1mm]\small Observo, escucho y pregunto.} &
\phasecard{2}{milcTurquesa}{milcGris}{Entender\\[-1mm]\small Sistematizo y ordeno.} &
\phasecard{3}{milcMagenta}{milcGris}{Hacer\\[-1mm]\small Construyo y aplico.} &
\phasecard{4}{milcVino}{milcGris}{Cierre\\[-1mm]\small Evalúo y reflexiono.}
\end{tabularx}


\puente{Empezamos viendo un correo real para descubrir sus partes.}

\milcestacion{milcVerde}{Estación 1 de 3 · Escucha}

\begin{softbox}{milcVerde}{milcGris}{Escena para escuchar}
\textbf{Actividad 1 · IDENTIFICA --- Detectar las partes del correo} (10 min · individual). Lee con calma el siguiente correo de ejemplo. Después \textbf{en el cuaderno} dibuja un cuadro y adentro escribe las \textbf{5 partes que descubriste}. No tienes que saber los nombres técnicos todavía; descríbelas con tus palabras: \emph{``La parte de arriba que dice de qué trata el correo''}, \emph{``El saludo''}, etc. \textbf{Ejemplo a leer:} \texttt{Para: lcardenas@sormariajuliana.edu.co} · \texttt{Asunto: Permiso para faltar el viernes 22} · ``Buenos días profe Álvaro. Le escribo porque el viernes 22 tengo cita médica a las 9 de la mañana y voy a faltar a su clase. Me comprometo a ponerme al día con las actividades de la semana. Por favor avíseme qué debo traer la próxima clase. Muchas gracias por su atención. Atentamente, Lara Cárdenas Martínez, grado 6B.''
\end{softbox}

{\sffamily\bfseries\fontsize{8}{10}\selectfont\textcolor{milcNegro!55}{MARCA CUANDO LO HAYAS HECHO}}
\milccheckrow{Leí el correo de ejemplo completo.}
\milccheckrow{Identifiqué 5 partes diferentes.}
\milccheckrow{Las describí con mis palabras (no con palabras técnicas).}

\begin{infoband}{milcVerde}


\textbf{Lo que va al cuaderno (Actividad 1 · IDENTIFICA).} Título: «Actividad 1 --- Las partes del correo». \textbf{Formato:} 5 frases numeradas, una por parte que descubriste. \textbf{Extensión:} 5 líneas, una por parte. \textbf{Sabes que terminaste cuando} tienes 5 cosas escritas y sabes señalarlas en el correo de ejemplo.
\end{infoband}

\puente{Después de ver el ejemplo, vas a aprender las 5 partes con sus reglas exactas.}

\milcestacion{milcTurquesa}{Estación 2 de 3 · Entender}

\begin{softbox}{milcTurquesa}{milcGris}{Lo que vas a entender}
Lo que descubriste solo es la base. Ahora vas a aprender los \textbf{nombres técnicos} de esas 5 partes y \textbf{una regla} para cada una. Son las mismas 5 partes que se usan en cualquier correo formal, desde un colegio hasta una empresa.
\end{softbox}

\milcpaso{1}{milcTurquesa}{\textbf{Parte 1 --- Destinatario (``Para:'').} El correo de la persona que va a recibir. Regla: \textbf{revísalo dos veces} antes de enviar. Un punto o letra de más y el correo se va a un desconocido. Ejemplo: \texttt{lcardenas@sormariajuliana.edu.co}, no \texttt{lcardenas@sormariajuliana.com}.}
\milcpaso{2}{milcTurquesa}{\textbf{Parte 2 --- Asunto.} Una frase corta que dice de qué trata el correo. Regla: \textbf{escribe el asunto pensando en el que recibe}, no en ti. Bien: \emph{``Permiso para faltar el viernes 22''}. Mal: \emph{``Hola profe''} o \emph{``Es urgente''}.}
\milcpaso{3}{milcTurquesa}{\textbf{Parte 3 --- Saludo formal.} Cómo empiezas. Regla: \textbf{usa el nombre y el cargo}. Bien: \emph{``Buenos días profe Álvaro''}, \emph{``Estimada señora rectora''}. Mal: \emph{``Qué más''}, \emph{``Hola''} a secas.}
\milcpaso{4}{milcTurquesa}{\textbf{Parte 4 --- Cuerpo.} Lo que quieres decir. Regla: \textbf{tres bloques cortos}. Bloque 1: qué pides o cuentas. Bloque 2: por qué. Bloque 3: qué pides que pase ahora. No más de 5 líneas en total. Nadie quiere leer un correo larguísimo.}
\milcpaso{5}{milcTurquesa}{\textbf{Parte 5 --- Despedida y firma.} Cómo cierras. Regla: \textbf{despedida cortés + tu nombre completo + tu curso}. Bien: \emph{``Muchas gracias por su atención. Atentamente, Lara Cárdenas Martínez, grado 6B''}. Mal: \emph{``Chao''} sin nombre.}

\begin{drawbox}{milcTurquesa}{Las 5 partes del correo formal}
\textbf{1. Destinatario} (``Para:'') --- a quién va. \\[1mm]
\textbf{2. Asunto} --- de qué trata, en una línea. \\[1mm]
\textbf{3. Saludo formal} --- nombre + cargo. \\[1mm]
\textbf{4. Cuerpo} --- 3 bloques cortos, máximo 5 líneas. \\[1mm]
\textbf{5. Despedida y firma} --- cortesía + tu nombre completo + curso.
\end{drawbox}

\begin{softbox}{milcMostaza}{milcGris}{Errores comunes y su corrección}
\textbf{Errores que se repiten cuando empiezas a escribir correos.} (1) Asunto vacío o solo ``Hola'' --- el que recibe no sabe de qué se trata. (2) Saludo informal a un adulto que no conoces --- \emph{``Hey profe''} no abre puertas. (3) Cuerpo enorme con muchas explicaciones --- la persona ocupada no termina de leer. (4) Firmar solo con apodo o nombre incompleto --- la persona no sabe quién escribió. (5) Escribir \textbf{TODO EN MAYÚSCULAS} --- en internet eso es como gritar.
\end{softbox}

\begin{infoband}{milcTurquesa}


\textbf{Lo que va al cuaderno (Actividad 2 · EXPLICA).} Título: «Actividad 2 --- Las 5 partes con sus reglas». \textbf{Formato:} tabla de 2 columnas. Columna 1: nombre de la parte. Columna 2: la regla con tus palabras (no copiada). \textbf{Extensión:} 5 filas (una por parte). \textbf{Sabes que terminaste cuando} podrías corregir el correo de un compañero solo mirando tu tabla.
\end{infoband}

\puente{Con las 5 partes claras, activas tu correo y escribes los 3 correos formales en el cuaderno.}

\milcestacion{milcMagenta}{Estación 3 de 3 · Hacer}

\begin{softbox}{milcMagenta}{milcGris}{Cómo vas a construir tu producto}
\textbf{Actividad 3 · APLICA --- Tu correo activo + 3 correos formales} (35 min · individual). Vas a hacer 2 cosas. \textbf{Primero:} activar tu correo institucional en el computador (\textbf{Paso A}). \textbf{Segundo:} escribir 3 correos formales en el cuaderno (\textbf{Paso B}). El profe te entrega tu usuario y contraseña inicial. Sigue los pasos en orden y con calma.
\end{softbox}

\milcpaso{1}{milcMagenta}{\textbf{Paso A.1 --- Abre el navegador.} En el computador abre Chrome o Edge. En la barra de arriba escribe \texttt{gmail.com} y pulsa Enter. Si te pide cuenta, escribe el correo que te dio el profe: \texttt{tunombre@sormariajuliana.edu.co}.}
\milcpaso{2}{milcMagenta}{\textbf{Paso A.2 --- Ingresa la contraseña inicial.} El profe te dio una contraseña temporal (ejemplo: \texttt{Sormaria2026}). Escríbela exactamente igual, respetando mayúsculas y números. Si te equivocas, pulsa \emph{``Olvidé mi contraseña''} y pide ayuda al profe.}
\milcpaso{3}{milcMagenta}{\textbf{Paso A.3 --- Cambia la contraseña inicial.} El sistema te pedirá una contraseña nueva. \textbf{Reglas de contraseña fuerte:} mínimo 8 caracteres, mezcla mayúsculas y minúsculas, agrega 2 números, agrega 1 signo (\texttt{!} o \texttt{?}). Ejemplo válido: \texttt{Mango2026!}. Ejemplo malo: \texttt{lara123}. \textbf{Escríbela en el cuaderno en una hoja aparte y no la muestres a nadie}.}
\milcpaso{4}{milcMagenta}{\textbf{Paso B.1 --- Correo 1 al profe de Tecnología.} En el cuaderno, escribe en formato de correo (las 5 partes) un mensaje al profe Álvaro pidiéndole que te recuerde la actividad de la próxima clase. Asunto sugerido: \emph{``Pregunta sobre la actividad de la próxima clase''}.}
\milcpaso{5}{milcMagenta}{\textbf{Paso B.2 --- Correo 2 al coordinador.} Escribe un correo al coordinador del colegio (\emph{``Buenos días señor coordinador''}) explicándole por qué crees que la sala de tecnología debería tener algo (un ventilador, mejor luz, etc.). Asunto sugerido: \emph{``Sugerencia para la sala de Tecnología''}.}
\milcpaso{6}{milcMagenta}{\textbf{Paso B.3 --- Correo 3 a un familiar lejano.} Escribe un correo a un tío o abuela que vive lejos contándole 3 cosas que estás aprendiendo este año. Asunto: \emph{``Contándote cómo voy en el colegio''}. Saludo: \emph{``Querida tía Marta''} (formal pero familiar).}

\begin{drawbox}{milcMagenta}{Antes de mostrar al profe}
\checkbox Mi correo institucional abre sin problemas. \\[1mm]
\checkbox Cambié la contraseña por una fuerte (la anoté en hoja aparte). \\[1mm]
\checkbox Los 3 correos están escritos en el cuaderno. \\[1mm]
\checkbox Cada correo tiene las 5 partes (destinatario, asunto, saludo, cuerpo, firma). \\[1mm]
\checkbox El asunto de cada correo dice de qué se trata sin abrir el correo. \\[1mm]
\checkbox Firmé con nombre completo y grado.
\end{drawbox}

\begin{softbox}{milcMagenta}{milcGris}{Plantilla de trabajo}
\textbf{Plantilla del correo formal en el cuaderno.} \textit{Para:} (correo de la persona). \textit{Asunto:} (frase corta). \textit{Saludo:} ``Buenos días/tardes [nombre + cargo]''. \textit{Cuerpo:} 3 bloques cortos --- qué pasa, por qué, qué pides. \textit{Despedida:} ``Muchas gracias por su atención. Atentamente, [nombre completo], grado [tu curso]''.
\end{softbox}

\begin{infoband}{milcMagenta}


\textbf{Lo que va al cuaderno (Actividad 3 · APLICA).} Título: «Actividad 3 --- Mi correo activo + 3 correos formales». \textbf{Formato:} en hoja aparte (privada): contraseña. En hoja del cuaderno: los 3 correos uno debajo del otro, con las 5 partes claras. \textbf{Extensión:} 1 página y media. \textbf{Sabes que terminaste cuando} cualquier compañero puede leer tu correo y entender exactamente qué pides.
\end{infoband}

\puente{El producto es doble: el correo configurado en el computador y los 3 correos escritos.}

\milcestacion{milcMostaza}{Producto de la sesión}
\textbf{Correo institucional activo + 3 correos formales en el cuaderno}.
\begin{softbox}{milcMostaza}{milcGris}{Criterios de calidad}
(1) El correo institucional abre y la contraseña nueva es fuerte (8+ caracteres, mayúsculas, números, signo). (2) Los 3 correos están escritos completos en el cuaderno. (3) Cada uno tiene las 5 partes exactas (destinatario, asunto, saludo, cuerpo, despedida). (4) Los asuntos son frases que cuentan de qué trata, no genéricos. (5) El cuerpo de cada correo cabe en 5 líneas o menos.
\end{softbox}

\puente{Antes de salir, miras tus 3 correos y verificas que cada uno tiene las 5 partes completas.}

\milcestacion{milcVino}{Evaluación liberadora · cinco dimensiones}

\begin{softbox}{milcVino}{milcGris}{Sentido de la evaluación}
Evaluar no es solo revisar una nota. Es reconocer qué aprendí, qué cambió en mí, qué emociones manejé, cómo me relaciono con la ciudad y la comunidad, y qué le dejaré a quien viene detrás.
\end{softbox}

\textbf{Mi reflexión liberadora en cinco dimensiones.} Completa cada frase iniciada:

\small
\begin{tabularx}{\textwidth}{>{\bfseries\RaggedRight\arraybackslash}p{3.4cm}Y>{\RaggedRight\arraybackslash}p{4.6cm}}
\rowcolor{milcVino}\color{white}Dimensión & \color{white}\bfseries Mi respuesta (frame starter) & \color{white}\bfseries Algo que descubrí\\
1. Personal & Lo que más cambió en mí fue \linead{6.5cm} & \linead{4.2cm}\\
2. Emocional & La emoción más fuerte fue \linead{2.5cm} y la regulé \linead{3cm} & \linead{4.2cm}\\
3. Ciudadana & En democracia esto importa porque \linead{6cm} & \linead{4.2cm}\\
4. Local (Cartago) & En mi barrio sería útil para \linead{6.5cm} & \linead{4.2cm}\\
5. Intergeneracional & A mi abuela/abuelo le diría \linead{6.5cm} & \linead{4.2cm}\\
\end{tabularx}
\normalsize

\begin{infoband}{milcVino}
\textbf{Para la próxima sustentación.} Vuelve a esta tabla en seis meses. Verás que las respuestas cambian — no porque lo aprendido cambie, sino porque tú cambias. La evaluación liberadora es una conversación contigo a través del tiempo.
\end{infoband}


\puente{Cerramos con 3 voces que te ayudan a entender por qué la palabra escrita pesa más de lo que crees.}

\milcestacion{milcGradeDark}{Triángulo de pensamiento · cierre filosófico}

\begin{softbox}{milcVino}{milcGris}{Dussel · lente del nosotros}
\textit{Escribir con cortesía no es agacharse: es respetar al otro y respetarte a ti mismo.}

{\footnotesize\itshape\color{milcNegro!70}--- formulación de esta guía, al modo de Enrique Dussel. No es una cita textual.}

Cuando escribes un correo formal, no estás siendo \emph{``ñoño''} ni perdiendo tu estilo. Estás haciendo lo contrario: estás diciendo \textbf{``yo merezco respeto y te respeto a ti''}. Por eso saludas, por eso firmas. Cuanto más serio el asunto, más cuidadoso el correo. Esa es la regla en cualquier lugar del mundo.

\textbf{Cuando le escribo a un adulto, ¿le escribo como si yo mismo mereciera que me lean en serio?}
\end{softbox}
\linead{\linewidth}\\[1mm]
\linead{\linewidth}

\begin{softbox}{milcMostaza}{milcGris}{Estoicismo · Séneca · cuidado interior}
\textit{Antes de enviar una carta, léela tres veces. La primera para corregir, la segunda para acortar, la tercera para asegurarte de que dice lo que querías decir.}

{\footnotesize\itshape\color{milcNegro!70}--- formulación de esta guía, al modo de Séneca. No es una cita textual.}

Hace 2000 años, en Roma, ya se sabía esto: las cartas (hoy correos) se mandan después de leerse varias veces. Antes de pulsar ``enviar'', léelo. ¿Tiene las 5 partes? ¿Cabe en 5 líneas? ¿Suena cortés? \textbf{El que envía sin releer manda errores en serio.} El que relee, manda mensajes que funcionan.

\textbf{¿Estoy enviando con prisa, o estoy revisando lo que digo antes de mandarlo?}
\end{softbox}
\linead{\linewidth}\\[1mm]
\linead{\linewidth}

\begin{softbox}{milcTurquesa}{milcGris}{Floridi · lente de la infoesfera}
\textit{Cada correo es una pequeña promesa: prometes que lo que dices es verdad y que estarás del otro lado para responder.}

{\footnotesize\itshape\color{milcNegro!70}--- formulación de esta guía, al modo de Luciano Floridi. No es una cita textual.}

Cuando mandas un correo institucional, no es como un mensaje de WhatsApp que se borra. \textbf{Queda guardado, queda firmado, queda contigo.} Si pides algo, lo que dices te compromete. Por eso el correo institucional es serio: lo que escribes ahí es como tu palabra escrita.

\textbf{Lo que estoy escribiendo en este correo, ¿lo voy a sostener si me preguntan después?}
\end{softbox}
\linead{\linewidth}\\[1mm]
\linead{\linewidth}

\begin{infoband}{milcNegro}
\textbf{Nota para el docente.} Las tres formulaciones son del autor de la guía, no citas textuales de los pensadores: recogen su lente para pensar el tema, no sus palabras. Si vas a citarlos en clase, remite a la obra original.
\end{infoband}

\milcestacion{milcVino}{Mi autoevaluación · marca una casilla por fila}

{\small
\setlength{\extrarowheight}{2.2pt}
\begin{tabularx}{\textwidth}{>{\RaggedRight\arraybackslash\bfseries}p{3.0cm}YYY}
\rowcolor{milcVino}\color{white}\bfseries Criterio & \color{white}\bfseries Lo logré & \color{white}\bfseries En proceso & \color{white}\bfseries Necesito apoyo\\[.6mm]
Comprensión del tema & \milcopcion{Explico las ideas con mis palabras.} & \milcopcion{Reconozco las ideas, falta precisión.} & \milcopcion{Necesito repasar lo básico.}\\[.8mm]
\rowcolor{milcGris}Pensamiento computacional & \milcopcion{Usé los pilares al armar mi producto.} & \milcopcion{Usé algunos pilares.} & \milcopcion{Necesito releer la estación 2.}\\[.8mm]
Aplicación y producto & \milcopcion{Mi producto cumple los criterios.} & \milcopcion{Está armado, con ajustes pendientes.} & \milcopcion{Aún está en construcción.}\\[.8mm]
\rowcolor{milcGris}Reflexión 5 dimensiones & \milcopcion{Respondí cada una con honestidad.} & \milcopcion{Respondí algunas con detalle.} & \milcopcion{Mis respuestas son muy generales.}\\[.8mm]
Triángulo de pensamiento & \milcopcion{Conecté las 3 voces con mi trabajo.} & \milcopcion{Conecté al menos una.} & \milcopcion{No las relaciono con mi proceso.}\\[.8mm]
\end{tabularx}}

\vspace{3mm}
\textbf{Hoy entendí que:}\\[1mm]
\linead{\linewidth}\\[1mm]
\linead{\linewidth}

\textbf{Mi compromiso para la próxima sesión es:}\\[1mm]
\linead{\linewidth}\\[1mm]
\linead{\linewidth}

\begin{softbox}{milcMostaza}{milcGris}{Cita de despedida · Marco Aurelio}
\textit{``El obstáculo en el camino se vuelve el camino. Lo que estorba al camino se vuelve camino.''}\\[1mm]
Cada error de hoy, bien observado, se convierte en pieza del camino que pisarás mañana.
\end{softbox}

\small
\begin{tabularx}{\textwidth}{>{\bfseries}p{3.6cm}Y}
\rowcolor{milcGradeDark!12}Nombre completo: & \linead{10cm}\\
Fecha: & \linead{4cm} \quad Curso/grupo: \linead{4cm}\\
Mi firma: & \linead{9cm}\\
\end{tabularx}
\normalsize

\begin{infoband}{milcGradeDark}
\textbf{Para la próxima sesión.} Llega con tu correo activo y los 3 correos escritos en el cuaderno. La próxima clase vas a conocer la \textbf{netiqueta}: 10 acuerdos básicos para comportarte con respeto en internet. Lo que aprendiste hoy ya cuenta como netiqueta sin que te dieras cuenta: saludo, cortesía, cuidar lo que dices. Ahora vas a verlo completo.
\end{infoband}

\end{document}
